% Updated October 26, 2026 at 4:43 PM EDT (America/New_York)
\documentclass[14pt,aspectratio=169]{beamer}
% Use the same MTU theme and course identification as the ODE lectures.
\ifdefined\pdfmapfile\pdfmapfile{+cm-super-t1.map}\fi
\usepackage{beamerthemeMTU}
\usepackage{amsmath,amssymb}
\course{MA3521 --- Introduction to ODEs}
\author{Allan Struthers}
\institute{Math Sciences}
\title{Int by Parts Review}
\lecturelabel{Review}
\date{}
\setbeamersize{text margin left=10mm,text margin right=10mm}

\begin{document}

\begin{frame}{Int by Parts Review: Backwards product rule!}
\begin{itemize}
\item Integrate and rearrange the  product rule to get
\[
d(uv)=u\,dv+v\,du 
\Longrightarrow
\boxed{\int u\,dv=uv-\int v\,du} + C
\]
Differentiate $u$, integrate $dv$ include $+C$ in the final integral.
\item Make $\int v\,du$ simpler. Useful choices 
that simplify when differentiated are $u=t$, $u=\ln(t)$, $u=\arctan(t)$, \dots

\begin{itemize}
\item Identify $u$ and $dv$.
\item Differentiate $u$ to get $du$. Integrate $dv$ to get $v$.
\item Substitute into $\int u\,dv=uv-\int v\,du$.
\end{itemize}
\end{itemize}
\end{frame}

\begin{frame}{Worked example \quad $\displaystyle\int \ln(t)\,dt$}
For $t>0$, write the integrand as $\ln(t)\cdot 1$.
\[
\boxed{\begin{aligned}
u&=\ln(t)\\
du&=\frac{1}{t}\,dt
\end{aligned}}
\qquad
\boxed{\begin{aligned}
dv&=dt\\
v&=t
\end{aligned}}
\]
\[
\int\ln(t)\,dt
=t\ln(t)-\int t\,\frac{1}{t}\,dt
=t\ln(t)-\int 1\,dt
\]
\[
\boxed{\displaystyle\int\ln(t)\,dt=t\ln(t)-t+C}
\]
\end{frame}

\begin{frame}{Worked example \quad $\displaystyle\int t e^{2t}\,dt$}
	\vspace{-2ex}
\[
\boxed{\begin{aligned}
u&=t\\
du&=dt
\end{aligned}}
\qquad
\boxed{\begin{aligned}
dv&=e^{2t}\,dt\\
v&=\frac{e^{2t}}{2}
\end{aligned}}
\]
\[
\int t e^{2t}\,dt
=\frac{t e^{2t}}{2}-\int\frac{e^{2t}}{2}\,dt
=\frac{t e^{2t}}{2}-\frac{1}{2}\int e^{2t}\,dt
\]
\[
\boxed{\displaystyle\int t e^{2t}\,dt
=\frac{t e^{2t}}{2}-\frac{e^{2t}}{4}+C}
\]

Remember the $1/2$ each time you integrate $e^{2t}$.

\end{frame}







\end{document}


